\section{Evaluation Framework: Security and Cost}
\label{sec:framework}

% UNIFIED framework section (grill #5): ONE framing for how we compare schemes, split into
% Security (categorical taxonomy) and Cost (quantitative perf+fidelity). The categorical/
% quantitative split is stated ONCE here; nothing downstream re-derives it. Side-channel
% finding lives in §\ref{sec:synthesis}; the validation case study in \cref{app:validation}.

We compare confidential Transformer-inference and RAG schemes along two kinds of
dimension: \emph{security}---categorical properties of who the adversary is and what
stays hidden (\cref{sec:threat})---and \emph{cost}---the quantitative price of that
protection, in performance and fidelity (\cref{sec:comparison}). This single split runs
through the paper: the security dimensions yield the threat-model-fit criterion of our
deployment-readiness lens (\cref{sec:intro}). We also fix here how schemes are
\emph{grouped} for comparison---by use case, not by family.

\subsection{Taxonomy: Threat Model and Security}
\label{sec:threat}

Confidential Transformer-inference and RAG systems are not built against a single
adversary. Each scheme makes its own assumptions about who the attacker is, what it
knows, what it can observe, and what must stay hidden---and those assumptions, not the
mechanisms alone, decide whether a defence is adequate. We classify the security of each
scheme along five \emph{categorical} axes (\cref{sec:threat-axes}), place a
representative set of schemes on them (\cref{sec:threat-landscape}), and treat the
attacker's \emph{capability budget} as a cross-cutting concern (\cref{sec:threat-budget}).
A field-wide side-channel blind spot emerges from the exercise (\cref{ssec:sidechannel}),
and the one empirical case study in this SoK (\cref{app:validation}) instantiates a single
point in this space.

\subsubsection{Axes}
\label{sec:threat-axes}

We classify each scheme along five axes.

\paragraph{1. Adversary behavior.}
\emph{Honest-but-curious} (HbC): the operator runs the protocol faithfully but
mines every observable for private information. \emph{Malicious/active}: the
operator may additionally deviate---drop, replay, reorder, or tamper with
computation. The obfuscation and TEE-split schemes we study---AloePri
\citep{lin_aloepri}, GELO \citep{belikov_gelo}, TwinShield
\citep{xue_twinshield}, ObfuscaTune \citep{frikha_obfuscatune}---adopt the HbC
model standard in cloud settings. Opal \citep{kaviani_opal} is the outlier: it
defends a malicious host that may tamper with everything outside the enclave.
Behavior decides whether confidentiality alone suffices or integrity machinery
is also required.

\paragraph{2. Trust anchor (TCB).}
What the scheme roots its guarantees in. We distinguish: \emph{none} (purely
cryptographic or obfuscation-based---AloePri assumes no trusted hardware); a
\emph{CPU-TEE} (Intel SGX/TDX, e.g.\ TwinShield's root of trust); a
\emph{GPU-TEE} (a confidential H100/H200/B200-class accelerator, e.g.\ GELO's
trusted device); and \emph{multi-party non-collusion} (the assumption
underlying MPC schemes such as PermLLM \citep{permllm}). A larger TCB buys
performance or simplicity at the price of a stronger trust assumption.

These anchors collapse onto a three-level \emph{hardware-trust} ladder that we use to
organize the systematization and its scope (\cref{sec:intro}):
\begin{description}
  \item[Tier 1.] \emph{No trusted hardware}---confidentiality rests on cryptography or
        statistics alone (MPC/FHE, static obfuscation, DP); trust anchor \emph{none}.
        ``None'' is a statement about what must be \emph{trusted not to leak}---no trusted
        third party or hardware---not about where the \emph{work} runs: several Tier-1
        schemes are not thin-client. Client-side-PIR/ORAM retrieval (Pacmann, Compass) and
        2PC inference where the user is a compute party (PermLLM) make the client
        \emph{load-bearing}---it stores access state and runs the oblivious-access or
        traversal loop, delegating only bulk storage to the server. We mark these
        ``none; client-relied'' and report the resulting client burden, a first-order
        deployment factor, in \cref{sec:synthesis}.
  \item[Tier 2.] \emph{CPU-TEE only}---a commodity CPU enclave (Intel SGX/TDX, AMD
        SEV-SNP) is trusted; any GPU used for acceleration is \emph{untrusted}.
  \item[Tier 3.] \emph{Confidential accelerator}---a confidential GPU (NVIDIA H100/B200
        CC) or other accelerator-TEE is trusted end-to-end.
\end{description}
The Tier-2$\to$Tier-3 jump is the sharpest trust escalation on this axis---trusting the
accelerator silicon end-to-end---so a scheme's tier is the single most consequential
entry in its row. The hybrid split's defining move (\cref{sec:hybrid}) is to deliver
Tier-3-class GPU performance while remaining at Tier~2.

\paragraph{3. Attacker knowledge.}
\emph{Black-box} (queries only), \emph{white-box on weights} (knows the public
model---the realistic case for open-weight Transformer models), or \emph{white-box plus the
obfuscation scheme} (knows the algorithm, ignorant only of an ephemeral
per-session secret). We adopt Kerckhoffs's principle throughout: across the
corpus the adversary knows the model and the method and is denied only the
fresh secret---AloePri's token permutation, GELO's per-batch mixing matrix.
This is the honest assumption for open-weight deployments, and precisely the
regime in which static schemes are most stressed: STIP \citep{stip} and
related static-permutation designs are secure only while the weights are
private, and collapse once an attacker who knows the public weights $W$ can
solve for the secret transform from the obfuscated $W'$ \citep{belikov_gelo}.

\paragraph{4. Confidentiality scope.}
What the scheme keeps hidden, and what it concedes---two faces of one question.
The protected \emph{asset} may be the input prompt; the input \emph{and} output
(AloePri protects both); the model weights as well (ObfuscaTune and TwinShield
additionally shield a proprietary model from the cloud); or the access pattern
over a corpus (Opal and the RAG-storage schemes). Its dual is the \emph{residual
leakage}---what the adversary still observes despite the scheme: obfuscated
service I/O (AloePri), mixed activations in untrusted VRAM (GELO, ObfuscaTune,
the TEE-split designs), bounded storage access patterns (RemoteRAG
\citep{cheng2025remoterag}), or cosine-distance structure (Caprise's
DCPE \citep{ye_caprise}). The access-pattern-hiding schemes (Opal, Pacmann,
Compass) sit on the protection side of this dual, not the leakage side: ORAM and
PIR are precisely what reduce their residual leak to an operation type or
nothing at all. Microarchitectural, cache, timing, and power
side channels are a residual-leakage category that \emph{every} surveyed scheme
places out of scope---a uniformity we elevate to a finding
(\cref{ssec:sidechannel}) rather than carry as a per-scheme cell. Two schemes
that agree on every other axis can diverge sharply here, and as the validation case study
(\cref{app:validation}) shows, it is the residual leakage---not the stated
protection goal---that an attack ultimately exploits.

\paragraph{5. Security basis.}
What the confidentiality guarantee \emph{rests on}---the form of the argument and
the assumption behind it---independent of \emph{which} entity is trusted
(axis~2). We distinguish: \emph{cryptographic} (provable under a
computational-hardness assumption, as in FHE and MPC under LWE/OT);
\emph{information-theoretic} (no computational assumption, as in one-time
masking); \emph{differential privacy} (a formal $(\varepsilon,\delta)$ or metric
bound); \emph{hardware-enforced} (security reduces to TEE isolation and
attestation correctness); and \emph{heuristic} (no formal guarantee---resistance
is shown experimentally, as in most static obfuscation). This axis separates
schemes the others cannot: an FHE scheme and a static-obfuscation scheme can
share an identical row on axes~1--4---trust anchor \emph{none},
honest-but-curious, white-box, protecting the input---yet the former offers
cryptographic indistinguishability while the latter offers only an empirically
measured resistance that the case study (\cref{app:validation}) shows can be far
weaker than its authors report. Security basis is thus the sharpest predictor of
whether a confidentiality claim survives scrutiny, and it is the axis the
deployment-readiness criterion of \emph{threat-model fit} (\cref{sec:intro})
reads most heavily.

\subsubsection{The Landscape}
\label{sec:threat-landscape}

\cref{tab:threat-taxonomy} places a representative scheme from each of the five
mechanism families on the five axes; the full per-family rosters appear in
\cref{sec:crypto}--\cref{sec:hybrid}. Four patterns stand out. First, the
obfuscation family clusters tightly at \emph{HbC + white-box}: these schemes
explicitly assume the attacker knows the open-source model and the method, and
stake everything on an ephemeral secret. Second, the trust anchor varies more
than the adversary behavior does: schemes that face nearly the same attacker still
root trust in very different places---nothing, a CPU-TEE, or a GPU-TEE---which is the real
axis of disagreement in the field. Third, only Opal and Compass
defend a \emph{malicious} operator---Opal a host that may tamper with computation outside
the enclave, Compass a retrieval server that may deviate from the agreed search protocol
(caught by its commitment and integrity checks, with no trusted hardware); every other
scheme assumes the operator follows the protocol. Fourth---and most consequential---the \emph{security basis} cleaves the
table where no other axis can: the obfuscation rows rest on heuristic,
experimentally-measured resistance, while PermLLM and the cryptographic-retrieval
schemes rest on hardness assumptions and the TEE designs on hardware. Several
heuristic and cryptographic schemes alike anchor trust in \emph{nothing}
(trust anchor \emph{none}), so this distinction is invisible without the basis
axis---yet it is exactly what the case study (\cref{app:validation}) turns on.

\begin{table*}[t]
  \centering
  \footnotesize
  \caption{Threat models across confidential Transformer-inference and RAG
    schemes, on the five axes of \cref{sec:threat-axes}. \emph{Knows}:
    wb=white-box weights, +s=+scheme. \emph{Protects} and \emph{Residual leak}
    are the two facets of confidentiality scope (asset hidden vs.\ what the server
    still sees): act=activations (VRAM), ap=access patterns, ct+dist=ciphertext +
    cosine-distance structure (the distance, not the ciphertext, is the leak). \emph{Basis}=security basis (crypto=cryptographic,
    DP=differential privacy, HW=hardware-enforced, heur.=heuristic;
    *=cryptographic but distance-leaking). Side channels are out of scope for
    every row (\cref{ssec:sidechannel}).}
  \label{tab:threat-taxonomy}
  \begin{tabular}{@{}lllllll@{}}
    \toprule
    Scheme & Behavior & Trust anchor & Knows & Protects & Residual leak & Basis \\
    \midrule
    \multicolumn{7}{@{}l}{\textit{Static obfuscation}}\\
    AloePri      & HbC        & none           & wb+s & in+out      & obf.\ I/O+wts & heur. \\
    STIP         & HbC        & none (3-party) & wb   & input       & activations   & heur. \\
    \midrule
    \multicolumn{7}{@{}l}{\textit{TEE / hybrid}}\\
    GELO         & HbC        & GPU-TEE        & wb   & input       & activations   & heur.+HW \\
    ObfuscaTune  & HbC        & CPU/GPU-TEE    & wb   & data+model  & activations   & heur.+HW \\
    TwinShield   & HbC+verif. & CPU-TEE + SS   & wb   & data+model  & activations   & crypto+HW \\
    Opal         & malicious  & TDX + GPU-TEE  & ---  & data+access & op-type only  & crypto+HW \\
    PipeLLM      & HbC        & GPU-TEE        & ---  & data+model  & PCIe act.     & HW \\
    \midrule
    \multicolumn{7}{@{}l}{\textit{Cryptographic / DP / RAG}}\\
    PermLLM      & HbC        & 2PC+dealer     & wb   & input        & shares       & crypto+heur. \\
    RemoteRAG    & HbC        & none           & ---  & query        & query+ap     & DP \\
    Caprise      & HbC        & none           & ---  & doc+query    & ct+dist      & crypto*+DP \\
    Pacmann      & HbC        & none; clt.-relied & ---  & query+access & none (pub.\ DB) & crypto \\
    Compass      & malicious  & none; clt.-relied & --- & query+access & op-type only  & crypto \\
    \bottomrule
  \end{tabular}
\end{table*}

\needswork{rows are grounded via the EdgeQuake-indexed source PDFs (AloePri, GELO,
ObfuscaTune, TwinShield, Opal, PermLLM, RemoteRAG, Caprise, Pacmann, Compass);
PipeLLM/STIP from their preprints --- re-confirm at submission}

\subsubsection{The Attacker's Capability Budget}
\label{sec:threat-budget}

The five axes say what an adversary \emph{is}; they do not say how \emph{much}
it can spend. Yet the capability budget---compute, and especially the
observation horizon (a single inference vs.\ many adaptive runs)---is decisive,
because the surveyed schemes claim security against markedly different budgets.
AloePri is explicit that it ``aims to provide adjustable security for
constrained attackers in real-world scenarios, rather than offering ideal-world
security guarantees against worst-case attackers'' \citep{lin_aloepri}; its
protection rests on per-inference noise calibrated to a Rényi-metric DP budget.
GELO's guarantee is narrower still: per-batch non-identifiability under a fresh,
never-reused mixing matrix, so the attacker faces a \emph{single-batch} blind
source separation problem and---by GELO's own analysis---gains nothing across
batches \citep{belikov_gelo}. A worst-case adversary with unbounded compute and
unlimited adaptive observation defeats every \emph{static}-obfuscation scheme here
outright; a per-batch design like GELO still concedes nothing \emph{across} batches,
leaving only a single-batch blind-source-separation problem (hard, but not
information-theoretically precluded---and SGT's stronger information-theoretic claim is
itself contested, \cref{sec:obfuscation}). An evaluation that assumed such an adversary
would therefore measure nothing interesting and would strawman schemes that never
claimed that bar.

This is why \emph{threat-model fit}---not threat-model maximality---is one of the three
deployment-readiness criteria. A scheme that defends an adversary
stronger than a deployment faces (an unbounded, fully-adaptive attacker, or
microarchitectural side-channels most operators accept or mitigate out-of-band) pays for
that protection in performance and may price itself out of use; a scheme that defends
too weak an adversary fails to protect. The deployable point is the one whose adversary
is \emph{calibrated} to the deployment---an argument that right-sizes security rather
than relaxing it, and that we apply when scoring families in \cref{sec:synthesis}.

We treat the budget as an explicit, per-scheme parameter rather than a hidden
constant. Crucially, the ``what breaks'' case study of \cref{app:validation} attacks each
scheme \emph{strictly within its own constrained-attacker regime}: the attacker is white-box
on weights and method, is denied the per-session secret, and operates under the same compute
and observation horizon the scheme's own evaluation permits. A break under those conditions is
a break by the victim's own standard, not an artefact of an inflated adversary.

\subsection{Cost: Performance and Fidelity}
\label{sec:comparison}

The other two deployment-readiness criteria are quantitative: \emph{performance}
(overhead and latency) and \emph{fidelity} (impact on output quality). Comparing them
across papers is hard to do fairly: reported numbers come from different hardware, model
sizes, datasets, and metrics, so a raw side-by-side can mislead. We fix here a common set of
\emph{reference metrics} for both, and state how we group schemes when we compare them.

\paragraph{Normalizing performance.}
Wherever a paper reports its own plaintext baseline, we express cost as an
\emph{overhead multiplier} ($\times$ vs.\ plaintext on the same hardware and model),
which cancels the hardware and model to first order and is the most portable figure
across papers. Where only absolute numbers exist, we compare only \emph{within a
regime}---a (model-size class $\times$ hardware class) bucket (e.g.\ BERT-base encoder
on a single GPU; a 7B decoder over WAN; ANN at a fixed corpus size)---never across
regimes. Every cross-scheme figure here is \emph{indicative} by default: each is drawn from its own
paper's testbed---different hardware, model, and dataset---so it conveys an order of magnitude,
not a precise ratio. The one exception is a figure we mark \emph{co-measured}: two or more
schemes benchmarked \emph{together, in one paper on one testbed}, where a direct comparison is
sound. We flag co-measured figures explicitly; every unflagged figure is indicative.
Even \emph{within} a regime, three further conditions must match before a comparison is fair,
and we hold them fixed or flag a mismatch: the \emph{number of computing parties} (a 2PC, 3PC,
or dealer-aided protocol carries different cost structure), the \emph{numerical regime}
(quantized vs.\ full-precision schemes are not directly comparable), and the \emph{network
setting}---Chen et al.~\citep{chen_sok_pti} re-benchmark crypto-inference systems and find that
several tuned for one bandwidth/latency profile do not reach their reported performance under a
different one, so a WAN figure and a LAN figure are not interchangeable inside the same
model-and-hardware bucket.

\paragraph{One-hop relative normalization.}
Cryptographic-inference schemes overwhelmingly report cost \emph{relative to a prior secure
scheme}, not to plaintext. Where a scheme $A$ reports a $k\times$ factor against a baseline
$B$, and $B$ reports an absolute figure in the \emph{same} regime, we recover $A$'s absolute
cost as ${\approx}\,\mathrm{cost}(B)/k$ and flag it \emph{extrapolated} (1-hop via $B$); we
never chain beyond one hop, nor extrapolate across regimes. Because these baseline chains
terminate at other \emph{secure} schemes rather than plaintext, a vs-plaintext multiplier for
cryptographic inference is reported only as an order-of-magnitude estimate against a stated
plaintext reference for the same regime, and flagged as such.

\paragraph{Normalizing fidelity.}
Fidelity is the impact on output quality relative to the \emph{plaintext} model. Schemes
that compute the true function---cryptography, TEE, and invertible-mask hybrid
splits---are \emph{exact} (zero quality loss) and we mark them so. Approximate schemes
(obfuscation, DP, polynomial-approximation MPC) are reported as a \emph{delta} against
the plaintext baseline in the quality metric appropriate to the use case
(\cref{tab:yardstick}), so a recall drop in retrieval and an accuracy drop in generation
are never conflated.

\paragraph{Grouping: by use case, not by family.}
We \emph{systematize} the field by mechanism family (\cref{sec:crypto}--\cref{sec:hybrid})
because that is how the design space is structured. But we \emph{compare} performance and
fidelity \textbf{by use case}, because that is how a deployment decision is actually
made: an operator who needs private \emph{generation} weighs every mechanism that serves
generation---MPC, TEE, obfuscation, hybrid split---against one another, not the FHE
family against the obfuscation family in the abstract. Grouping by use case puts the real
alternatives side by side. We use four deployment use cases, aligned with the pipeline
stages of the gap-map (\cref{sec:synthesis}): \emph{generation} (decoder inference),
\emph{embedding} (encoder inference), and \emph{retrieval}, which we split by index
structure into \emph{vector retrieval} (flat/IVF nearest-neighbour search and vector
storage) and \emph{graph retrieval} (graph-based ANN and graph-semantic search/storage),
since the two impose different costs and threat surfaces. Each scheme is grouped by the
model class it actually evaluates---an encoder-only scheme is an embedding scheme, a
decoder scheme a generation one. \emph{Reranking} (cross-encoder scoring) is a pipeline
stage we track in the gap-map, but no surveyed scheme reports a standalone
private-reranking measurement, so we treat it as an open gap (\cref{sec:synthesis}) rather
than a cost bucket. \cref{tab:yardstick} fixes the canonical performance metric, fidelity
metric, and plaintext baseline for each.

\begin{table}[t]
  \centering \footnotesize
  \caption{Reference metrics: per-use-case canonical metrics against which
    \cref{sec:crypto}--\cref{sec:hybrid} report and \cref{sec:synthesis} compares.}
  \label{tab:yardstick}
  \begin{tabular}{@{}llll@{}}
    \toprule
    Use case & Performance metric & Fidelity metric & Plaintext baseline \\
    \midrule
    Generation & latency/token; $\times$ overhead & task acc.\ / perplexity / BLEU $\Delta$ & plaintext decode \\
    Embedding (encoder) & latency/seq; $\times$ overhead   & downstream task-acc.\ $\Delta$  & plaintext encode \\
    Vector retrieval & latency/query @corpus-size; comm & recall@$k$ $\Delta$               & plaintext flat/IVF ANN \\
    Graph retrieval  & latency/query; comm              & recall@$k$ / answer-acc.\ $\Delta$ & plaintext graph search/ANN \\
    \bottomrule
  \end{tabular}
\end{table}
\needswork{once the per-family numbers are filled, confirm each is normalized to the
matching row of \cref{tab:yardstick} and flagged co-measured vs indicative}
